\sectionkicker{Theme synthesis}
\subsection*{横断比較 — Agent Execution Stackはどこへ責務を移したか}
\addcontentsline{toc}{subsection}{横断比較 — Agent Execution Stackはどこへ責務を移したか}
1月のagent関連更新を同じ製品カテゴリとして並べるより、base modelの外側へ移った責務で分解すると変化が見えやすい。harness、context、action surface、access control、human controlは別々の実装面で拡張されていた。
\medskip
\begin{center}
\small
\renewcommand{\arraystretch}{1.16}
\begin{tabularx}{\linewidth}{@{}p{0.20\linewidth}X@{}}
\toprule
観察軸 & 1月の一次資料・論文から読めること \\
\midrule
harness / orchestration & OpenAIのCodex loopではuser、model inference、tools、conversation history、context-window managementをharness logicが編成する。model単体ではなく実行loop全体が設計対象になる。 \cite{src-fa5410010e0a,src-ae10eb4597fe} \\
context / permission & OpenAIのin-house data agentではmetadata、human annotation、code enrichment、institutional knowledge、memory、runtime contextを層として重ね、既存権限をpass-throughする。Codex loopとは別systemであり、組織内case studyとして読む。 \cite{src-fa5410010e0a,src-ae10eb4597fe} \\
desktop action surface & AnthropicのCoworkは1月12日時点のresearch previewとしてClaudeのagentic capabilityをdesktop workflowへ持ち込んだ。後日のavailabilityをこの時点へ逆投影しない。 \cite{src-f6aa679babd7} \\
behavior guidance & Anthropicが1月22日に公開したClaude constitutionはtraining processの一部としてbehaviorを方向付ける規範層であり、desktop execution surfaceとは役割が異なる。実出力が常に理想どおりになる保証でもない。 \cite{src-f6aa679babd7} \\
computer use & Gemini API changelogは1月29日、gemini-3-pro-previewとgemini-3-flash-previewにComputer Use tool supportを記録した。一次事実はpreview identifierへのAPI availabilityであり、behavior品質の独立評価ではない。 \cite{src-a24fd97eb3d0} \\
input / lifecycle & 同じGemini API changelogは1月8日のCloud Storage・public/private pre-signed database URL入力対応と1月12日のmodel lifecycle機能も記録する。action surfaceだけでなく入力経路とlifecycle管理も同じAPI面で拡張された。 \cite{src-a24fd97eb3d0} \\
terminal control loop & Mistral Vibe 2.0はterminal-native coding agentへcustom subagents、multi-choice clarification、slash-command skills、unified agent modesを加えた。1月のeventは基盤model familyの新設ではなくagent/harness側の更新である。 \cite{src-d8b7c1c36830} \\
distribution / access & Vibe 2.0はLe Chat Pro/TeamとPAYGまたはBYOK経路で提供されたとされる。実行harnessの機能だけでなく、どのcredential・product経路で使うかもagent deploymentの境界になる一方、生産性評価はvendor claimのままである。 \cite{src-d8b7c1c36830} \\
\bottomrule
\end{tabularx}
\renewcommand{\arraystretch}{1.0}
\end{center}
\normalsize
\begin{claimboundary}[この比較の境界]
この表は本号で既に検証・参照している一次資料と論文を横断比較の形へ再配置したものである。vendor / project / author が報告した性能値や優位性は独立再現済みの一般則へ変換せず、本文とSource Notesの帰属境界をそのまま維持する。
\end{claimboundary}
